\chapter{Sheaves, Field Coherence, and Global Consistency}

\section{Local Truth, Global Consistency}

Every chapter since Chapter 12 has quietly restricted its attention to a single point, a single trajectory, or a single region considered in isolation. Writing, projection, commit, fixed points, energy descent, and Hamiltonian oscillation were all stated as though the courtyard's various points had nothing to do with one another beyond sharing the same field. This was never quite true, and this chapter is where the omission is repaired. A field defined over a domain X is not merely a collection of independent stories, one per point; it is a single object, and something has to guarantee that what is locally admissible at each point actually agrees with its neighbors where their local pictures overlap. A wall that is locally coherent on its left half and locally coherent on its right half is not thereby coherent as a whole, unless the two halves also agree where they meet.

\section{The Field-Coherence Sheaf}

This is a local-to-global consistency question, and Chapter 10 already introduced the mathematical machinery suited to answering questions of this shape, applied there to cues rather than to the field itself. The same machinery applies here, over a different base. Let 𝓜 be the field's configuration space, as developed in Chapter 18, and let the embedded structure of the field's domain, written G ↪ 𝓜 following the embedding this book's supporting technical material already develops in full, provide a covering: a way of decomposing X into overlapping patches, each small enough to be checked for local admissibility on its own.

The field-coherence sheaf, \(F_{\mathrm{field}}\), assigns to each patch of this cover the set of locally admissible field configurations on that patch, admissible in the sense Chapter 18's C already made precise, and assigns to each inclusion of a smaller patch into a larger one the restriction of a configuration from the larger patch to the smaller. The gluing question this sheaf poses is the one section 21.1 raised informally: given locally admissible configurations on each patch of a cover, agreeing wherever two patches overlap, do they assemble into a single globally admissible configuration on the whole domain? When they do, \(F_{\mathrm{field}}\) is said to glue on that cover, and the resulting global section is exactly what this book has meant, since Chapter 6, by a coherent world.

\section{Hallucination as a Cohomological Obstruction}

It is worth being precise about what happens when gluing fails, because this book's supporting technical material has, since its earliest chapters, described hallucination as a cohomological obstruction without yet earning the right to use that language. The failure of local data to assemble into a global section is measured, in the branch of mathematics this construction borrows from, by a cohomology class: informally, a record of exactly how the overlaps disagree, which cannot itself be resolved by any choice of global configuration because the disagreement is already baked into the local data before any assembly is attempted. Write \(H^1\)(𝒢∞) for this obstruction, evaluated over the full cover 𝒢∞ implied by the field's embedding. A field state for which \(H^1\)(𝒢∞) = 0 is one whose local patches, wherever checked, agree on their overlaps and can be assembled without contradiction. A field state for which \(H^1\)(𝒢∞) ≠ 0 is one where no such assembly exists: two locally admissible, individually well-formed pieces of the field that simply cannot both be true of the same underlying world at once, a courtyard whose left half implies a doorway the right half's own local structure has no room for.

This gives hallucination a sharper meaning than this book has previously had available. It is not merely a proposal that violates a locally checkable rule, the kind of failure Chapter 14's projection already catches. It is a global failure that no amount of local checking, patch by patch, will ever surface, because each patch, examined on its own, may be perfectly admissible. Only the overlap reveals the problem, and only cohomology, rather than any pointwise constraint from Chapters 8 through 20, is equipped to detect it.

\section{From Cue Sheaves to World Sheaves}

Chapter 10 promised that its cue-coherence sheaf, \(F_{\mathrm{cue}}\) : \(\mathrm{CueCat}^{\mathrm{op}}\) → Set, would eventually be connected to whatever field-coherence construction this book developed, rather than left to stand as an unrelated structure sharing only a family resemblance. That promise is kept here. The connection runs through observation, in exactly the sense Chapter 6 first introduced it and Chapter 9 gave concrete form to. A cue, recall, is a signal arising from the field, and this is not a coincidence to be papered over with a freestanding functor invented for this chapter's convenience; it is \(O_\theta\), the observation operator this book has developed since Chapter 6, doing exactly the work this bridge requires. Write O for the map \(O_\theta\) induces from CueCat into Patch(𝓜), sending each cue's context to the patch of the field that observation actually samples to produce it. This is not a new primitive. It is \(O_\theta\), restricted to the role of relating cue-contexts to field-patches rather than fields to observations directly, and the direction this book adopts, cues arising downstream of patches rather than patches arising downstream of cues, follows from Chapter 9's own account of where a candidate ẑ comes from: an agent is cued by what it has already, at least partially, observed.

With O in hand, the bridge is a natural transformation,

\[ \eta : F_{\mathrm{cue}} \Rightarrow O^*F_{\mathrm{field}}, \]

where O*\(F_{\mathrm{field}}\) is the pullback of the field-coherence sheaf along O, assigning to each cue-context whatever local field-admissibility O maps that context onto. η states the compatibility condition this book has been assuming informally since Chapter 10 without proving: an affordance activated at a cue is only genuinely available if its realization, once selected by Chapter 13's policy and proposed as a \(\Delta\psi\), corresponds to an admissible transformation of the field patch that cue actually arose from. A ladder that cues climbing but whose realization would require a field transformation outside \(C_S\) ∩ \(C_z\) ∩ \(C_R\) is not, on this account, a live affordance at all, however strongly the cue has activated it, and η is precisely the map that catches this: an affordance in the image of η is guaranteed field-realizable; an affordance outside its image is a cue-level illusion, coherent from the agent's side and unsupported from the world's.

\section{Completing the Constraint Manifold}

Chapter 18 left C's definition with an acknowledged placeholder, \(C_A\), promised but not yet formalized. It can now be stated precisely: \(C_A\) is the set of field states for which every affordance activated by any cue arising from that state lies in the image of η, that is, every locally available action possibility corresponds to a genuinely admissible field transformation rather than a cue-level illusion of one. With \(C_A\) in place, C = \(C_S\) ∩ \(C_z\) ∩ \(C_R\) ∩ \(C_A\) is no longer missing a named member, though this book has been careful, since Chapter 18, not to claim the list is closed; Chapters 22 and 23 will each have occasion to ask whether further conditions belong in the intersection before this book is willing to call C complete.

\section{Which Maps Preserve Admissibility}

Chapter 18 also deferred a question about domain transport: when a topology repair or identity restructuring changes the field's domain from X to some new X′, which maps φ relating the two are actually entitled to carry C's hard constraints across, and which are not. The sheaf-theoretic apparatus this chapter has developed answers this precisely. φ is eligible exactly when it is a morphism of sites: when it sends covers of X to covers of X′ compatibly with restriction, so that \(F_{\mathrm{field}}\)'s gluing condition on the old domain transports to a well-posed gluing condition on the new one rather than to a collection of patches that no longer cohere into any sensible covering at all. A φ that scrambles the cover, merging patches that should remain distinct or splitting one patch into pieces with no coherent restriction maps between them, is not eligible, however smooth or structure-preserving it may appear from outside the sheaf's own bookkeeping, and any constraint transported along such a φ would arrive already broken. This is the condition Chapter 18 promised would be supplied here, and it was never separable from the sheaf construction this chapter exists to develop; transport and gluing were always the same requirement viewed from two directions.

\section{Looking Ahead}

This chapter has given the field's local-to-global consistency a precise treatment, named hallucination as the specific cohomological failure it always informally was, completed the bridge between cue-level and field-level coherence that Chapter 10 promised, and supplied C with its fourth named member. What it has not asked is what, if anything, stays exactly conserved as the field evolves under all of this machinery at once, beyond the energy this book's variational chapters have already discussed. Identity, residue, and the coherence a committed configuration carries forward all seem, informally, like they ought to be conserved in some sense stronger than merely persisting at a fixed point. Chapter 22 asks what conservation actually means here, and what, precisely, this framework is entitled to claim never simply disappears.
